\section{Mean-zero primitives and harmonic coefficients}\label{tsc:primitives}
For $r\ge1$, let $U_{m,r}$ be the unique mean-zero periodic distribution
satisfying $\partial_x^r U_{m,r}=C_m$. Equivalently,
\[
 \widehat U_{m,r}(n)=(2\pi\ii n)^{-r}\widehat C_m(n),\quad n\ne0,
 \qquad \widehat U_{m,r}(0)=0.
\]
These are represented by ordinary locally integrable functions. We give
the exact normalization because a primitive fixed at a base point is not
in general the same function.

For $N\ge0$, define complete homogeneous harmonic coefficients by
\begin{equation}\label{tsc:homogeneous}
 \sum_{j\ge0}h_j^{(N)}v^j
     =\prod_{k=1}^N(1-v/k)^{-1}
     =\exp\left(\sum_{k\ge1}\frac{H_N^{(k)}}k v^k\right).
\end{equation}
Thus $h_0^{(N)}=1$, $h_1^{(N)}=H_N$,
$h_2^{(N)}=(H_N^2+H_N^{(2)})/2$, and
\[
 h_3^{(N)}=\frac{H_N^3+3H_NH_N^{(2)}+2H_N^{(3)}}6.
\]
For $N=0$, all positive-index coefficients are zero.

\begin{proposition}[Explicit normalized primitive]\label{tsc:primitive-formula}
For $m\ge0$, $r\ge1$, and $0<x<1$,
\begin{align}\label{tsc:Ucoef}
 U_{m,r}(x)
 &=m![u^{m+1}]\frac{\Gamma(1+u)}{\Gamma(r+u)}\zeta(1-r-u,x)\notag\\
 &=\frac{(-1)^{m+1}m!}{(r-1)!}
       \sum_{j=0}^{m+1}\frac{h_{m+1-j}^{(r-1)}}{j!}
                   \zeta^{[j]}(1-r,x).
\end{align}
In particular,
\begin{equation}\label{tsc:Uone}
 U_{m,1}(x)=\frac{(-1)^{m+1}}{m+1}\zeta^{[m+1]}(0,x),
 \qquad U_{0,1}(x)=\frac\ell2-\log\Gamma(x).
\end{equation}
\end{proposition}
\begin{proof}
The meromorphic family $\Gamma(u)\calB_{r+u}$ has Fourier coefficients
$\Gamma(u)(2\pi\ii n)^{-r-u}$. Its coefficient $m![u^m]$ is therefore
the stated mean-zero primitive. On the open interval, use
\eqref{tsc:Bhurwitz} and move the factor $1/u$ to the coefficient index;
this proves the first equality in \eqref{tsc:Ucoef}. Next,
\[
 \frac{\Gamma(1+u)}{\Gamma(r+u)}
   =\frac1{(r-1)!}\prod_{k=1}^{r-1}(1+u/k)^{-1}
   =\frac1{(r-1)!}\sum_{j\ge0}(-1)^jh_j^{(r-1)}u^j.
\]
Multiplication by the Taylor expansion of $\zeta(1-r-u,x)$ gives the second
equality. The case $r=1$ and Lerch's formula give \eqref{tsc:Uone}.
Mean zero and the distributional derivative relation follow from the Fourier
construction; they are not extra constants to be chosen afterward.
\end{proof}
These one-function ladders belong to the classical Hurwitz primitive
calculus and its repository developments \cite{tsc:em2,tsc:shj}. Their
three-factor use is the next result.

\begin{theorem}[All primitive triple identities]\label{tsc:primitive-triple}
For $r_i\ge1$ and $m_i\ge0$,
\begin{align}\label{tsc:Utriple}
 \int_0^1\prod_{i=1}^3U_{m_i,r_i}(\{x-a_i\})\dd x
 ={}&(2\pi)^{-r_1-r_2-r_3}\left(\prod_{i=1}^3m_i!\right)\notag\\
 &\quad\cdot[\mathbf u^{\mathbf m+\mathbf1}]
     \left(\prod_{i=1}^3E(u_i)\right)
                         \calT(\mathbf r+\mathbf u;\mathbf a).
\end{align}
This is an ordinary convergent integral. It remains valid when some or all
shifts coincide, with the right side evaluated in its locally convergent
positive-order Tornheim representation.
\end{theorem}
\begin{proof}
For separated shifts, apply \eqref{tsc:master-eq} to
$\prod_i\Gamma(u_i)\calB_{r_i+u_i}$ and extract the three nonnegative
Laurent coefficients, exactly as in Theorem~\ref{tsc:stieltjes-triple}.
This gives \eqref{tsc:Utriple}.

Near a seam, $U_{m,1}$ has at worst logarithmic growth of finite degree.
For $r\ge2$ the possible singular term is a power $x^{r-1}$ multiplied
by a logarithmic polynomial, and the function remains bounded. Every
finite product in the theorem is therefore integrable, including at a
common seam. On the double-sum side, a neighborhood of $r_i\ge1$ is an
absolute-convergence domain even for trivial colors: logarithmic
spectral derivatives are dominated by a slightly decreased positive power,
and the worst central denominator is $mn(m+n)$. Summing first over
$m+n=N$ gives $2H_{N-1}/N^2$. This proves local normal convergence for
all the finitely many derivatives in question. Hence the separated formula
passes to coincident shifts on both sides. This uses no continuation of a
singular Stieltjes product to a collision.
\end{proof}

\begin{corollary}[Shifted centered log-Gamma cube]\label{tsc:loggamma-triple}
Let $G(x)=\log\Gamma(x)-\ell/2$. Then
\begin{equation}\label{tsc:Gtriple}
 \int_0^1\prod_{i=1}^3G(\{x-a_i\})\dd x
   =-\frac1{(2\pi)^3}
       \left.\mathscr D_A\calT(\mathbf s;\mathbf a)\right|_{\mathbf s=(1,1,1)}.
\end{equation}
\end{corollary}
At coincident shifts this is another representation of the known
Bailey--Borwein--Borwein cubic log-Gamma evaluation, not a new evaluation
of that old special case \cite{tsc:bbb,tsc:gamma-chapter}. Arbitrary shifted
mixed-index and mixed-primitive versions follow directly from
\eqref{tsc:Utriple} without another Fourier calculation.
